\section{Justificación}
\label{sec:justification}

El estudio y la gestión de los ecosistemas marinos frente al cambio climático global exigen un cambio de paradigma metodológico en la ingeniería ambiental regional. Tradicionalmente, los modelos de pronóstico a gran escala (como ROMS o Delft3D) han operado bajo un enfoque empírico de caja negra, aproximando las complejas interfaces bentónicas ---como las cubiertas coralinas y las matrices de sedimentos del Parque Nacional Natural Tayrona--- como fronteras planas, estáticas y uniformes. Para resolver el transporte de masa, estos modelos dependen de coeficientes de transferencia de masa y arrastre ajustados por calibración empírica (\emph{ad hoc}). Sin embargo, estas parametrizaciones clásicas carecen de validez física ante regímenes climáticos disruptivos y no lineales, limitando severamente la capacidad predictiva de institutos locales como el INVEMAR.

La verdadera innovación científica no surge de la repetición de metodologías tradicionales a macroescala, sino de operar de manera disruptiva en las interfaces interdisciplinarias. Como ingeniero químico enfocado en la física de fluidos compleja, esta propuesta rompe con el empirismo convencional al introducir un enfoque multiescala fundamentado en primeros principios y fenómenos de transporte avanzados. En lugar de formular suposiciones macroscópicas arbitrarias, este proyecto desciende a la microestructura del flujo ($Re \ll 1$) utilizando la Dinámica Stokesiana para modelar de forma determinista la cinemática discreta, las fuerzas de lubricación y las interacciones hidrodinámicas de muchos cuerpos en las fronteras vivas del ecosistema.

Al acoplar matemáticamente estas trayectorias microscópicas con la ecuación macroscópica de Advección-Difusión-Reacción (ADR), los parámetros macroestructurales críticos ---el Tensor de Dispersión Hidrodinámica Efectiva ($\bm{D}_{\text{eff}}$) y la Tasa de Captura Kinética ($R$)--- se derivan \emph{a priori} mediante rigurosos protocolos de escalamiento basados en la varianza asintótica a tiempos largos. Integrar la precisión computacional de la ingeniería química en la resolución de problemas de conservación costera no solo elimina la dependencia de coeficientes de sintonización empíricos, sino que transforma el modelado ambiental regional de un ejercicio de ajuste de datos a una ciencia predictiva rigurosa, dotando a la Universidad Nacional de Colombia de una línea de investigación de vanguardia global con un profundo valor territorial.
